\section{Background and related work}
\label{sec:background}

An OWL~2 ontology is a finite set of axioms interpreted under the OWL~2 direct
semantics \cite{w3c2012direct,w3c2012profiles}.  We write TBox for the
terminological axioms, ABox for assertions about individuals, and RBox for role
axioms.  Classification computes the entailed subsumption relation between
named classes, the unsatisfiable named classes, and the consistency of the
ontology; KM reports $A\sqsubseteq B$ exactly when every model of the ontology
interprets $A$ as a subset of $B$.  Realisation computes, for every named
individual, the named classes it belongs to in every model.

Tableau procedures decide a subsumption $A\sqsubseteq B$ by trying to construct
a model of the ontology in which an element belongs to $A$ and to the
complement of $B$; the ontology entails the subsumption exactly when every
attempt closes.  Blocking makes the construction finite, absorption rewrites
general axioms into triggered rules to reduce nondeterminism, and dependency
tracking lets the search backjump over irrelevant choices
\cite{horrocks1999optimizing,horrocks2007implementation,tsarkov2007optimizing}.
HermiT's hypertableau calculus reduces nondeterminism further by resolving
entire clause bodies at once \cite{motik2009hypertableau,glimm2014hermit}.
Konclude adds a saturation pass that decides most subsumption tests before any
tableau search, extensive caching of satisfiable node labels, and parallel
classification \cite{steigmiller2014konclude,steigmiller2014coupling}.  FaCT++
and its Java port JFact, RacerPro, and Pellet and its continuation Openllet
represent the earlier generation of tableau systems
\cite{tsarkov2006factpp,palmisano2018jfact,haarslev2012racerpro,sirin2007pellet,galigator2019openllet}.
These systems cover the full language, and their cost grows with the number
and size of the models an ontology forces.

Consequence-based procedures derive subsumptions without constructing models.
For \EL{}, completion rules over a small set of normal forms compute all
named subsumptions in polynomial time \cite{baader2005el}; CEL implemented
them first \cite{baader2006cel}, and ELK added indexing and concurrent
saturation that classify SNOMED~CT in seconds
\cite{kazakov2011concurrent,kazakov2014elk}.  Kazakov's procedure carried the
approach to Horn $\mathcal{SHIQ}$ \cite{kazakov2009horn}.  Beyond Horn
ontologies a consequence-based procedure has to reason about disjunctions.
The context calculus behind Sequoia keeps, for every symbolic element, a
context whose clauses are derived by ordered resolution and exchanged with
predecessor and successor contexts, and it has been extended to number
restrictions, inverse roles, and nominals
\cite{bate2016sequoia,bate2018consequence,tenacucala2018consequence,tenacucala2021payg}.
These procedures are goal-directed and reuse consequences across subsumption
tests; a context that accumulates many incomparable disjunctions can exhaust
memory before it saturates.

Because each family fails on ontologies that another handles, hybrid designs
send an ontology, or a part of it, to the procedure that fits.  MORe extracts
a module inside a tractable profile and sends the remainder to an expressive
reasoner \cite{armasromero2012more}; PAGOdA brackets the answer between a
lower and an upper bound computed by a tractable reasoner and calls an
expressive one for the gap \cite{zhou2015pagoda}; R2O2* learns to rank
external reasoners from ontology features \cite{kang2020r2o2}.  KM also routes
on source features.  Its routed procedures are internal to one executable and
share one clause representation, and each must return a complete answer or
decline (Section~\ref{sec:procedures}).

Reuse of existing implementations has precedents on which KM's port of
Konclude builds.  JFact is a Java port of FaCT++ \cite{palmisano2018jfact},
Whelk reimplements ELK's algorithm in Scala as an immutable data structure,
and whelk-rs reimplements Whelk in Rust \cite{balhoff2024whelk,whelkrs}.  The
Rust ecosystem for OWL also includes the Horned-OWL parsing library
\cite{lord2024horned}, which KM uses for OWL/XML and RDF/XML input, and
OxidOWL \cite{sieve2026oxidowlcodingagents,oxidowl}, a tableau and hypertableau
reasoner for OWL~2~DL built on Horned-OWL and, like KM, developed with coding
agents.  OxidOWL is the system closest to KM in language coverage and
implementation language, and Section~\ref{sec:eval-oxidowl} compares the two.

The correctness of an OWL reasoner's implementation is usually established by
test suites and by agreement with other systems.  Konclude's authors compared
its taxonomies with HermiT's and analysed the differences by hand
\cite{steigmiller2014konclude}, Sequoia's evaluation hashed taxonomies and
reported that each matched at least one other reasoner
\cite{tenacucala2019sequoia}, and the ORE~2015 organisers adjudicated results
by majority vote \cite{parsia2017ore}.  A justification-based comparison of
four reasoners on BioPortal ontologies found correctness bugs in mature
systems \cite{lee2015correctness}.  Two alternatives to agreement exist.  A
reasoner can be derived from a mechanised calculus
\cite{hidalgo2014verified,ileri2024vel}, or an optimised reasoner can emit
evidence that a small verified checker validates \cite{baader2020certify}.  KM
follows the second approach for its optimised procedures and the first for
reference procedures on fragments of the language
(Section~\ref{sec:certification}).
